\[ %%%%%%%%%%%%%%%%%%%%%%%%%%%%%%%%%%%%%%%%%%%%%%%%%%%%%%%%%%%%%%%%%%%%%%%%%%%%%%%% \subsection{Physics-Guided Message Passing Mechanism} \label{sec:physics_message_passing} %%%%%%%%%%%%%%%%%%%%%%%%%%%%%%%%%%%%%%%%%%%%%%%%%%%%%%%%%%%%%%%%%%%%%%%%%%%%%%%% The physics-guided message passing mechanism represents the central computational component of PHYSGEN. While conventional graph neural networks learn interaction patterns directly from data, PHYSGEN introduces physical guidance into the information propagation process itself. The fundamental assumption is that physical interactions between system components are not arbitrary statistical relationships but are governed by underlying principles, including locality, conservation laws, symmetry constraints, and domain-specific interaction rules. Therefore, message propagation should be dynamically regulated according to both learned representations and physical consistency requirements. %%%%%%%%%%%%%%%%%%%%%%%%%%%%%%%%%%%%%%%%%%%%%%%%%%%%%%%%%%%%%%%%%%%%%%%%%%%%%%%% \subsubsection{Conventional Graph Message Passing} %%%%%%%%%%%%%%%%%%%%%%%%%%%%%%%%%%%%%%%%%%%%%%%%%%%%%%%%%%%%%%%%%%%%%%%%%%%%%%%% A standard graph neural network updates node representations through iterative message aggregation: \begin{equation} m_i^t = \sum_{j\in\mathcal{N}(i)} M_{\theta}(h_i^t,h_j^t,e_{ij}), \end{equation} where $M_{\theta}$ represents a learnable message function and $\mathcal{N}(i)$ denotes the neighborhood of node $i$. The node update is then computed as: \begin{equation} h_i^{t+1} = U_{\theta}(h_i^t,m_i^t), \end{equation} where $U_{\theta}$ represents the node update function. Although this formulation provides powerful representation learning capabilities, the learned interactions are not explicitly constrained by physical principles. As a result, the model may discover correlations that improve training accuracy while violating physically meaningful relationships. %%%%%%%%%%%%%%%%%%%%%%%%%%%%%%%%%%%%%%%%%%%%%%%%%%%%%%%%%%%%%%%%%%%%%%%%%%%%%%%% \subsubsection{Physics-Guided Interaction Modulation} %%%%%%%%%%%%%%%%%%%%%%%%%%%%%%%%%%%%%%%%%%%%%%%%%%%%%%%%%%%%%%%%%%%%%%%%%%%%%%%% PHYSGEN extends conventional message passing by introducing a physics-guided interaction operator: \begin{equation} m_i^t = \sum_{j\in\mathcal{N}(i)} \alpha_{ij}^{phys} M_{\theta}(h_i^t,h_j^t,e_{ij}), \end{equation} where: \begin{equation} \alpha_{ij}^{phys} \in [0,1] \end{equation} represents a physics-guided modulation coefficient controlling the contribution of each interaction. The coefficient is computed as: \begin{equation} \alpha_{ij}^{phys} = \Psi_{phys} ( h_i^t, h_j^t, e_{ij}, P_{ij} ), \end{equation} where $P_{ij}$ contains physical descriptors of the interaction. These descriptors may include: \begin{itemize} \item spatial distance; \item relative velocity; \item material properties; \item conservation relationships; \item local physical constraints; \item interaction topology. \end{itemize} Therefore, information propagation becomes dependent not only on learned similarity but also on physical relevance. %%%%%%%%%%%%%%%%%%%%%%%%%%%%%%%%%%%%%%%%%%%%%%%%%%%%%%%%%%%%%%%%%%%%%%%%%%%%%%%% \subsubsection{Physics-Constrained Aggregation} %%%%%%%%%%%%%%%%%%%%%%%%%%%%%%%%%%%%%%%%%%%%%%%%%%%%%%%%%%%%%%%%%%%%%%%%%%%%%%%% A central feature of PHYSGEN is the introduction of physically constrained aggregation. For a node state update: \begin{equation} h_i^{t+1} = U_{\theta} ( h_i^t, m_i^t ), \end{equation} the aggregated message must satisfy physical admissibility conditions: \begin{equation} \mathcal{C} ( h_i^{t+1} ) \leq \epsilon, \end{equation} where $\mathcal{C}$ represents a physical consistency functional and $\epsilon$ denotes an admissible tolerance. Examples of such constraints include: \begin{itemize} \item conservation of mass; \item conservation of energy; \item bounded state evolution; \item symmetry preservation; \item physically valid parameter ranges. \end{itemize} Unlike conventional physics-informed approaches, these constraints influence the computation during message propagation rather than being applied only after prediction. %%%%%%%%%%%%%%%%%%%%%%%%%%%%%%%%%%%%%%%%%%%%%%%%%%%%%%%%%%%%%%%%%%%%%%%%%%%%%%%% \subsubsection{Conservation-Aware Message Correction} %%%%%%%%%%%%%%%%%%%%%%%%%%%%%%%%%%%%%%%%%%%%%%%%%%%%%%%%%%%%%%%%%%%%%%%%%%%%%%%% To further improve physical consistency, PHYSGEN introduces an optional correction mechanism: \begin{equation} \tilde{m}_i^t = m_i^t - \lambda_c \nabla_h \mathcal{C}(h_i^t), \end{equation} where $\lambda_c$ controls the strength of physical correction. This operation allows the network to reduce physically inconsistent information flows before they affect future state evolution. Consequently, the message passing process becomes: \begin{equation} h_i^{t+1} = U_{\theta} ( h_i^t, \tilde{m}_i^t ). \end{equation} %%%%%%%%%%%%%%%%%%%%%%%%%%%%%%%%%%%%%%%%%%%%%%%%%%%%%%%%%%%%%%%%%%%%%%%%%%%%%%%% \subsubsection{Multi-Hop Physics Propagation} %%%%%%%%%%%%%%%%%%%%%%%%%%%%%%%%%%%%%%%%%%%%%%%%%%%%%%%%%%%%%%%%%%%%%%%%%%%%%%%% Complex physical phenomena often involve interactions occurring at multiple spatial and temporal scales. PHYSGEN therefore applies repeated physics-guided message passing: \begin{equation} H^{(l+1)} = \Phi_{phys}^{(l)} ( H^{(l)},G_t ), \end{equation} where $l$ denotes the message-passing layer. Each layer represents an increasingly higher-order approximation of physical interaction propagation. The resulting latent representation captures both: \begin{itemize} \item local physical interactions; \item emergent global system behavior. \end{itemize} %%%%%%%%%%%%%%%%%%%%%%%%%%%%%%%%%%%%%%%%%%%%%%%%%%%%%%%%%%%%%%%%%%%%%%%%%%%%%%%% \subsubsection{Comparison with Attention-Based Graph Networks} %%%%%%%%%%%%%%%%%%%%%%%%%%%%%%%%%%%%%%%%%%%%%%%%%%%%%%%%%%%%%%%%%%%%%%%%%%%%%%%% Graph attention mechanisms calculate interaction weights based primarily on learned feature similarity: \begin{equation} \alpha_{ij} = Attention(h_i,h_j). \end{equation} While effective for many machine learning tasks, such attention coefficients do not necessarily correspond to physically meaningful interactions. PHYSGEN extends this concept by introducing physics-aware attention: \begin{equation} \alpha_{ij} = Attention_{phys} ( h_i,h_j,e_{ij},P_{ij} ). \end{equation} Therefore, interaction strength is determined by a combination of learned representations and physical constraints. %%%%%%%%%%%%%%%%%%%%%%%%%%%%%%%%%%%%%%%%%%%%%%%%%%%%%%%%%%%%%%%%%%%%%%%%%%%%%%%% \subsubsection{Summary} %%%%%%%%%%%%%%%%%%%%%%%%%%%%%%%%%%%%%%%%%%%%%%%%%%%%%%%%%%%%%%%%%%%%%%%%%%%%%%%% The physics-guided message passing mechanism transforms graph neural networks from purely statistical interaction learners into physically regulated computational systems. By embedding physical reasoning directly into information propagation, PHYSGEN enables the model to preserve structural properties of dynamical systems while maintaining the flexibility of neural representation learning. This mechanism provides the foundation for subsequent components of PHYSGEN, including latent physical dynamics evolution and adaptive stability control for long-horizon forecasting. \] 